\chapter{Number theory and integer roots}
\label{ch:numbertheory}

Include \file{biginteger/ops/NumberTheory.h} or \file{biginteger/ops/Operations.h}
for the APIs in this chapter. The prime tests give exact answers below
$2^{64}$ and probabilistic answers above that range.

\section{Greatest common divisors and B\'ezout coefficients}
\label{sec:gcd}

\code{Gcd(a, b)} returns the nonnegative greatest common divisor of signed inputs.
\code{ExtendedGcd(a, b)} also returns coefficients $x$ and $y$ satisfying
$a x + b y = \gcd(a,b)$. Both functions accept zero. For two zero inputs, the
divisor and both coefficients are zero.

One Euclidean step replaces $(a,b)$ by $(b,a-qb)$, where $q$ is the quotient
of $a$ divided by $b$. The common divisors stay the same. Extended GCD applies
the same step to the coefficients of the two remainders, so they remain linear
combinations of the original inputs.

\begin{example}
$240=5\cdot46+10$, then $46=4\cdot10+6$, $10=6+4$, and
$6=4+2$. Substituting back gives $2=-9\cdot240+47\cdot46$.
The coefficient pair need only satisfy the identity.
\end{example}

BigMath uses Lehmer's algorithm to batch steps when both remainders have
the same limb count and have at least two limbs. It predicts quotients from
a common leading prefix of at most 62 bits and bounds what the omitted
low limbs could contribute. A batch is used only when those bounds certify
at least two steps. A two-by-two matrix applies the batch to the full
remainders and, for extended GCD, to the coefficient pairs. An exact
division step handles cases where a batch cannot be certified. Plain GCD
finishes with a scalar remainder and a single-limb GCD when the second
operand fits in one limb.

\subsection{Recursive half-GCD}
\label{sec:halfgcd}

For large inputs, BigMath uses \href{https://doi.org/10.1090/S0025-5718-07-02017-0}{M\"oller's difference-based half-GCD}.
Subtracting $q$ times one entry from the other is a determinant-one
matrix operation. It has an integer inverse, so the two entries keep
the same common divisors. A product of these matrices represents a
whole reduction and can also act on the B\'ezout coefficients.

With $N$ the larger bit length, the recursive step works on the leading
bits and applies the resulting matrix to the full operands. For two
entries wider than $s=\floor{N/2}+1$ bits, it preserves that lower bound
and reduces their absolute difference to at most $s$ bits. The difference
test and exact correction steps control the effect of the omitted low
bits. Below or at 2\,048 bits, the recursion uses bounded Lehmer batches
and exact steps.

The shorter operand must have at least 65\,536 bits for \code{Gcd} or
16\,384 bits for \code{ExtendedGcd}, and its bit length must exceed
$\floor{N/2}+1$. A wider imbalance falls through to exact division.
Plain GCD applies one high-half reduction at a time. Extended GCD
composes a larger transform and applies it to both coefficient pairs.
The recursive algorithm costs $\bigO(\M(N)\log N)$, with $\M(N)$ the
cost of multiplying $N$-bit integers.

\begin{perfnote}
For balanced 262\,144-bit operands, paired single-threaded tests on
the M1 Max measured 2.93$\times$ faster GCD and 4.58$\times$ faster
extended GCD versus the Lehmer-based implementation.
\href{https://github.com/mmurshed/biginteger/blob/main/BENCHMARK.md\#recursive-half-gcd}{Measurement details}
\end{perfnote}

\section{Integer square root with a remainder}
\label{sec:sqrtrem}

\code{SqrtRem(n)} returns $(s,r)$ such that
$s=\floor{\sqrt n}$, $n=s^2+r$, and $0\le r\le 2s$.
It accepts zero and throws for negative inputs. For example, \code{SqrtRem(80)}
gives $(8,16)$ because $80=8^2+16$ and $9^2=81$.

For one or two limbs, a floating-point estimate supplies a starting root.
The code checks it with integer arithmetic and corrects it by integer Newton
iteration when necessary. Inputs too small for a recursive split likewise
use Newton iteration, seeded from leading bits.

For larger inputs, BigMath uses Zimmermann's divide-and-conquer square root.
Choose a limb-aligned radix $D$ and write $n=aD^2+bD+c$. Recursively compute
$h=\floor{\sqrt a}$ and $r_h=a-h^2$. Divide $r_h D+b$ by $2h$ to get a quotient
$q$ and remainder $t$. Expanding the trial root $hD+q$ gives
\[
n-(hD+q)^2=tD+c-q^2 .
\]
The split keeps $h\ge D/2$, so if the trial remainder is negative, one
decrement of the trial root suffices. To update the remainder after the
decrement, add twice the trial root minus one.

\begin{example}[ with radix 10]
For $n=2759$, the blocks are $27$, $5$, and $9$. The high root is $5$ with
remainder $2$. Dividing $25$ by $10$ gives $q=2$ and $t=5$, so the
root is $52$ and the remainder is $50+9-4=55$.
\end{example}

\newpage
\section{Integer roots of other degrees}
\label{sec:nthroot}

\code{RootRem(n, degree)} returns $(x,r)$ with $r=n-x^e$ for degree
$e$. The root is rounded toward zero. Negative inputs require an odd
degree. For a nonnegative input, $x^e\le n<(x+1)^e$. Degree zero throws,
degree one returns the input with remainder zero, and degree two uses
\code{SqrtRem}.

\begin{example}
\code{RootRem(-80, 3)} returns $(-4,-16)$ since $(-4)^3=-64$ and
$(-5)^3=-125$.
\end{example}

Magnitudes that fit in 64 bits use native binary search for other
degrees. A nonzero magnitude has root one if the degree is at least
its bit length. For the other inputs, BigMath first finds an exact leading
prefix of the root by binary search. To find the prefix with $s$ low
root bits removed, take the integer $e$th root of
$\floor{|n|/2^{es}}$. A prefix $p$ then bounds the full root from
above by $(p+1)2^s$.

BigMath refines the prefix in stages of increasing precision. Each
stage starts just above the root for its truncated input $A$ and
uses the integer Newton step
\[
  x_{\mathrm{next}} =
  \floor{\frac{(e-1)x+\floor{A/x^{e-1}}}{e}} .
\]
The first step that would stay at $x$ or increase it signals that
$x$ is the integer root for that stage. Exponentiation stops early
when a trial power exceeds $A$. The final stage uses the full input
and forms the remainder from the final power.

\begin{perfnote}
On the M1 Max, the sampled 64-bit \code{RootRem} call at degree three
took 0.094\,\textmu s in BigMath and 0.106\,\textmu s in GMP. At
262\,144 bits, it took 1.55\,ms in BigMath and 0.75\,ms in GMP.
\href{https://github.com/mmurshed/biginteger/blob/main/BENCHMARK.md\#integer-nth-root-and-remainder}{Measurements and degree sweep}
\end{perfnote}

\newpage
\section{Modular powers}
\label{sec:powmod}

\code{PowMod(a, e, m)} returns $a^e\bmod m$ in $[0,m)$.
The base may be signed, the exponent must be nonnegative, and the
modulus must be positive. An exponent of zero returns $1\bmod m$,
which is zero when $m=1$.

Repeated squaring processes the binary digits of $e$ and reduces each
intermediate product modulo $m$. For a dense exponent, a sliding window
uses a table of odd powers to handle several bits at once. Sparse or short
exponents use the binary path. In $7^{13}\bmod 11$, the powers for
exponents $1,4,8$ have residues $7,3,9$ and their product has residue $2$.

\begin{itemize}
  \item A single-limb modulus uses machine-word arithmetic.
  \item An odd modulus of at least 64 limbs and an exponent of at least
        64 bits use FastREDC by default. Otherwise, odd multi-limb moduli
        through 128 limbs use limb-by-limb Montgomery reduction for
        exponents of at least 16 bits.
  \item An even modulus of at most 128 limbs with a power-of-two factor
        $2^k$, where $0<k<\log_2 B$, can use the Chinese remainder
        theorem for exponents of at least 16 bits. The implementation
        splits off $2^k$ and applies the usual dispatch to the odd part.
  \item Other cases use the usual multiplication and division dispatchers.
\end{itemize}

The runtime depends on the operands and the exponent's bits.

\subsection{Montgomery reduction and FastREDC}
\label{sec:fastredc}

For an odd $n$-limb modulus $m$, let $R=B^n$. Montgomery form represents
$x$ by $xR\bmod m$. Multiplying two encoded residues gives $T<m^2$.
Computing $T/R\bmod m$ restores the encoding. Let $m'=-m^{-1}\bmod R$.
Then
\[
q=(T m')\bmod R, \qquad u=\frac{T+qm}{R}.
\]
The low $n$ limbs of $T+qm$ cancel. Dropping them gives $u<2m$, so one
conditional subtraction puts the result in $[0,m)$.

The base case cancels one limb at a time. FastREDC doubles the known
width of the inverse to precompute $m'$ once and uses a low product for
$q$. The initial product or square goes through BigMath's usual
dispatcher. With CRT enabled, reductions from 512 limbs can recover
the high part of $qm$ from a cyclic product. Write $qm=\ell+RH$, where
$\ell=(-T)\bmod R$ is already known. If $W=qm\bmod(B^L-1)$, subtract
$\ell$ and multiply by $B^{L-n}$ modulo $B^L-1$, a limb rotation, to recover $H$.
The bound $H<m$ makes that recovery unique. The cyclic path rounds $L$
up to a power of two within the standard CRT transform limit.

\begin{perfnote}
In paired single-threaded M1 Max runs, FastREDC sped up dense
8\,192-bit powering by 1.23$\times$. A 65\,536-bit modulus with a
257-bit exponent sped up by 1.57$\times$. These ratios include setup.
\href{https://github.com/mmurshed/biginteger/blob/main/BENCHMARK.md\#subquadratic-montgomery-reduction}{Measurement details}
\end{perfnote}

\newpage
\section{Primality testing and finding the next prime}
\label{sec:prime}

\code{IsProbablePrime(n, rounds)} rejects values below 2 and screens small
factors before running Miller--Rabin. The test writes $n-1=2^s d$ with $d$
odd. For a witness base $a$, it computes $a^d\bmod n$ and squares the residue
up to $s-1$ times. The first value may be $1$ or $-1$ modulo $n$. Otherwise,
one of the later squares must reach $-1$ for the witness to pass.

The deterministic word path uses three fixed bases through 32 bits and
seven through the rest of the unsigned 64-bit range. For larger values it
screens base 2, then samples uniformly from $[2,n-2]$ by rejection.
With $k$ independent uniform random witnesses, the chance that any
fixed odd composite passes all $k$ is at most $4^{-k}$. The default is
32 random witnesses. Both prime APIs require a positive round count.
Miller--Rabin uses the FastREDC cutoff with $d$ as the exponent and
reuses the precomputed inverse across witnesses in the same test.
The default build can test witnesses in parallel when the candidate
has at least 1,024 bits and at most 128 limbs and the caller requests
at least eight random rounds. Wider candidates can use the pool inside
individual products.

\code{NextProbablePrime(n, rounds)} searches strictly above $n$ and returns
2 for $n<2$. For example, a search starting at 29 returns 31. For
word-size starting values it tests odd candidates directly. For larger
starting values, it marks multiples of small odd primes in segments of
1,024 odd candidates. Each small prime's next offset carries forward to the
next segment. The sieve uses primes below 4,096 through 512 bits and below
32,768 above that size.

The first sieve survivor is tested on its own. In the parallel size range
(at least 1,024 bits and at most 128 limbs), the search can then screen
up to eight survivors at a time with base 2. It applies the randomized
tests to the passing candidates in ascending order, so completion order
cannot select a later result.

\begin{practice}
\file{BENCHMARK.md} records separate GMP comparisons for these operations,
including dense and sparse exponents and ordinary and long prime gaps.
In a Release build with GMP available, build \code{number\_theory\_vs\_gmp}.
Use \code{gcd-large} for recursive GCD or \code{powmod-large} for FastREDC.
The \code{gcd}, \code{sqrtrem}, \code{powmod}, \code{prime}, and
\code{nextprime} modes isolate the other measurements.
\end{practice}

\newpage
\section{Factorials and binomial coefficients}
\label{sec:combinatorics}

\code{Factorial(n)} computes $n!$, including $0!=1$. A native table covers
$0$ through $20$. For larger inputs, a product tree packs consecutive
factors into limbs and multiplies the packed groups in balanced pairs.
The tree gives the multiplication dispatcher similarly sized operands.

\subsection{Prime swing}

Prime swing starts at $n=512$ with 64-bit limbs and at $64$ with
32-bit limbs. Let $O(n)$ be the odd part of $n!$. The recursion uses
\[
  O(n)=O(\floor{n/2})^2
  \prod_{\substack{p\le n\\p\ \mathrm{odd\ prime}}}
       p^{\,\sum_{j\ge1}(\floor{n/p^j}\bmod 2)}.
\]
For each odd prime, the exponent counts the powers $p^j$ for which
$\floor{n/p^j}$ is odd. BigMath sieves once up to $n$ and reuses
prefixes of the prime list at smaller recursive sizes. The product
tree combines the swing factors. A final shift restores the
$n-\operatorname{popcount}(n)$ factors of two.

\begin{example}
For $10!$, the odd part of $5!$ is $15$ and the swing is $3^2\cdot7$.
Thus $O(10)=15^2\cdot63=14\,175$. Shifting by eight bits gives
$10!=3\,628\,800$.
\end{example}

\subsection{Binomial coefficients by exact division or prime powers}

\code{Binomial(n, k)} returns zero for $k>n$ and otherwise replaces $k$
by $\min(k,n-k)$. Through 32 terms, it packs batches of numerator
and denominator factors into limbs. Each batch ends at a binomial
coefficient, so \code{DivideExact} can divide the accumulated product
(\cref{sec:exactdivision}). Wider off-center cases build two packed
product trees and divide once.

For $n\ge512$ and reduced $k\ge\floor{n/8}$, the central path builds
the answer from prime powers. The exponent of $p$ is
$v_p(n!)-v_p(k!)-v_p((n-k)!)$, where
$v_p(j!)=\sum_{i\ge1}\floor{j/p^i}$. This path reuses the factorial
sieve and packed product tree, then restores the power of two.

\begin{perfnote}
On the M1 Max, the default 64-bit build took 28.01\,ms for
$1\,000\,000!$ and 5.61\,ms for $\binom{1\,000\,000}{500\,000}$.
GMP took 78.47\,ms and 5.91\,ms on those fixtures.
Use \code{number\_theory\_vs\_gmp factorial} or
\code{number\_theory\_vs\_gmp binomial} to reproduce the comparisons.
\end{perfnote}

\newpage
\section{Chinese remaindering with a reusable product tree}
\label{sec:crtplan}

\code{ChineseRemainderPlan} takes a nonempty list of pairwise coprime
word moduli $m_i\ge2$ and owns its precomputed data. Write
$M=\prod_i m_i$. \code{Reduce(value)} returns the least nonnegative
residue for each modulus in input order, including when the input is
negative. \code{Reconstruct(residues)} returns the unique integer in
$[0,M)$ with those residues.

\subsection{What the plan precomputes}

The plan packs consecutive moduli into groups whose products fit in
one word. A tree stores the product of each pair of child groups or
subtrees. The root holds $M$. A pass down the tree computes
$c_i=M/m_i\bmod m_i$ and precomputes its inverse $t_i$ modulo $m_i$.
If any inverse is missing, the moduli share a factor and construction
throws.

\subsection{Downward remainders and upward reconstruction}

For reduction, pass the input remainder down the tree, taking it
modulo each child's product. At a leaf, a word reduction supplies the
residue for the group. Each group then reduces to its individual
moduli. For reconstruction with residues $r_i$, the weighted terms
are
\[
  S=\sum_i \bigl(r_i t_i\bmod m_i\bigr)\frac{M}{m_i}.
\]
Modulo $m_i$, every term except the $i$th vanishes. At a tree node,
the implementation scales each child's contribution by the other
child's product and adds them. Leaves use word arithmetic and the
root reduces the result modulo $M$.

\begin{example}
With moduli $(3,5,7)$ and residues $(2,3,2)$, the cofactors are
$(35,21,15)$ and their inverses are $(2,1,1)$ in the respective
moduli. The weighted sum is $35+63+30=128$, so the representative
modulo $105$ is $23$.
\end{example}

\begin{practice}
Reuse one plan for repeated calls with the same moduli. The benchmark
times setup separately. For 8\,192 moduli of 61 bits on the M1 Max,
reconstruction took 3.17\,ms in the default 64-bit build and
3.77\,ms in FLINT. FLINT was faster at the smaller sampled sizes.
\href{https://github.com/mmurshed/biginteger/blob/main/BENCHMARK.md\#product-tree-chinese-remaindering}{Setup, reduction, and reconstruction timings}
\end{practice}

\newpage
\section{Bounded rational reconstruction}
\label{sec:rationalreconstruction}

\code{ReconstructRational(a, m, N, D)} looks for a reduced fraction
$n/d$ with $|n|\le N$, $0<d\le D$, and $n\equiv ad\pmod m$.
It accepts signed residues and normalizes them to $[0,m)$. An empty
optional means there is no fraction within the bounds. The call
requires $m>1$, $N\ge0$, $D>0$, and $2ND<m$.

The strict bound makes the fraction unique. If two candidates share
the residue, $m$ divides $n_1d_2-n_2d_1$. Its absolute value is at
most $2ND$, so it must be zero. The positive denominator and
reduced form fix both signs and the common scale.

Extended Euclid starts with $m$ and $a$ and tracks the coefficient of
$a$ in each remainder. It stops at the first remainder at most $N$.
That remainder and its coefficient give the numerator and denominator
after sign normalization. BigMath checks the denominator bound and
that the two are coprime. It uses Lehmer batches and recursive half-GCD
only while both remainders stay above the stopping bound
(\cref{sec:halfgcd}).

\begin{example}[ after Chinese remaindering]
The residues $(3,1,2)$ modulo $(5,7,11)$ give $a=288$ modulo $385$.
With $N=D=5$, the successive remainders are $97$, $94$, and $3$.
The coefficient of $288$ at remainder $3$ is $-4$, since
$3=3\cdot385-4\cdot288$. The normalized fraction is $-3/4$.
\end{example}

\begin{perfnote}
On the M1 Max, half-GCD reduced the two sampled 262\,144-bit
reconstructions from 77--119\,ms to about 21\,ms. FLINT took
11--14\,ms. With FLINT and GMP available, the
\code{rational\_reconstruction\_vs\_flint} target reproduces the comparison.
\href{https://github.com/mmurshed/biginteger/blob/main/BENCHMARK.md\#bounded-rational-reconstruction}{Measurements and test fixtures}
\end{perfnote}
